% This file is generated from the corresponding Obsidian note.
% Source: Teoricos/GS - Teo9.md
\chapter{Subvariedades}
\label{ch:teo-9}
\section{Subvariedades}
\begin{bookdefinition}{Inmersión}
Sean $M$ y $N$ variedades suaves y sea

\[
F:M\to N
\]
una función suave. Se dice que $F$ es una \textbf{inmersión} si para todo $p\in M$ la diferencial

\[
(dF)_p:T_pM\to T_{F(p)}N
\]
es inyectiva.

\end{bookdefinition}
\begin{bookdefinition}{Subvariedad}
Se dice que el par $(M,F)$ es una \textbf{subvariedad} de $N$ si

\[
F:M\to N
\]
es una inmersión y ademas es inyectiva.

\end{bookdefinition}
\begin{bookdefinition}{Subvariedad incrustada}
Se dice que $(M,F)$ es una \textbf{subvariedad incrustada} de $N$ si además

\[
F:M\to F(M)\subseteq N
\]
es homeomorfismo entre $M$ y $F(M)$ , donde $F(M)$ lleva la topología heredada de $N$ En ese caso se dice también se dice que $F$ es una \textbf{incrustación suave}.

\end{bookdefinition}
\section{Ejemplos}
\begin{bookexample}{Inclusion}
Supongamos que $m<n$ y definimos

\[
c:\mathbb R^m\to\mathbb R^n,\qquad c(x^1,\dots,x^m)=(x^1,\dots,x^m,0,\dots,0).
\]
Queremos ver que $c$ es una incrustación suave de $\mathbb R^m$ en $\mathbb R^n$.

\end{bookexample}
\begin{bookproof}{}
\begin{itemize}
\item \textbf{$c$ es suave}
\end{itemize}
\begin{enumerate}
\item Cada función coordenada de $c$ es suave:
\end{enumerate}
\[
c^i(x)=x^i\quad\text{si }1\le i\le m,
\]
y

\[
c^i(x)=0\quad\text{si }m+1\le i\le n.
\]
\begin{enumerate}[start=2]
\item Por lo tanto $c$ es $C^\infty$.
\end{enumerate}
\begin{itemize}
\item \textbf{$c$ es inyectiva}
\end{itemize}
\begin{enumerate}
\item Si
\end{enumerate}
\[
c(x^1,\dots,x^m)=c(y^1,\dots,y^m),
\]
entonces, comparando las primeras $m$ coordenadas,

\[
x^i=y^i,\qquad i=1,\dots,m.
\]
Luego

\[
(x^1,\dots,x^m)=(y^1,\dots,y^m).
\]
Por lo tanto $c$ es inyectiva.

\begin{itemize}
\item \textbf{Cálculo del diferencial}
\end{itemize}
\begin{enumerate}
\item Usamos cartas identidad:
\end{enumerate}
\[
\varphi=\operatorname{id}_{\mathbb R^m},\qquad \psi=\operatorname{id}_{\mathbb R^n}.
\]
\begin{enumerate}[start=2]
\item Entonces
\end{enumerate}
\[
\widehat c=\psi\circ c\circ\varphi^{-1}=c.
\]
\begin{enumerate}[start=3]
\item Y con sus respectivas bases sabemos
\end{enumerate}
\[
[dc_{x}]_{\mathcal{B_{1}}}^{\mathcal{B_{2}}}=J\widehat c(\varphi(x))=Jc(x)=\begin{pmatrix}1&0&\cdots&0\\0&1&\cdots&0\\\vdots&\vdots&\ddots&\vdots\\0&0&\cdots&1\\0&0&\cdots&0\\\vdots&\vdots&\ddots&\vdots\\0&0&\cdots&0\end{pmatrix}.
\]
\begin{enumerate}[start=4]
\item Así,
\end{enumerate}
\[
dc_x(v^1,\dots,v^m)=(v^1,\dots,v^m,0,\dots,0).
\]
\begin{enumerate}[start=5]
\item Si $dc_x(v)=0,$ entonces
\end{enumerate}
\[
(v^1,\dots,v^m,0,\dots,0)=(0,\dots,0),
\]
y por lo tanto

\[
v^1=\cdots=v^m=0.
\]
\begin{enumerate}[start=6]
\item Luego $v=0$, y entonces
\end{enumerate}
\[
dc_x:T_x\mathbb R^m\to T_{c(x)}\mathbb R^n
\]
es inyectiva.

\begin{itemize}
\item \textbf{$c$ es homeomorfismo sobre su imagen}
\end{itemize}
\begin{enumerate}
\item La imagen de $c$ es
\end{enumerate}
\[
c(\mathbb R^m)=\mathbb R^m\times\{0\}^{n-m}\subseteq\mathbb R^n.
\]
\begin{enumerate}[start=2]
\item Definimos
\end{enumerate}
\[
\pi:c(\mathbb R^m)\to\mathbb R^m,\qquad \pi(y^1,\dots,y^n)=(y^1,\dots,y^m).
\]
\begin{enumerate}[start=3]
\item Esta aplicación es continua porque es la restricción de la proyección continua
\end{enumerate}
\[
\operatorname{pr}:\mathbb R^n\to\mathbb R^m,\qquad \operatorname{pr}(y^1,\dots,y^n)=(y^1,\dots,y^m).
\]
\begin{enumerate}[start=4]
\item Además,
\end{enumerate}
\[
\pi(c(x^1,\dots,x^m))=(x^1,\dots,x^m),
\]
es decir,

\[
\pi\circ c=\operatorname{id}_{\mathbb R^m}.
\]
\begin{enumerate}[start=5]
\item Y si $y\in c(\mathbb R^m)$, entonces
\end{enumerate}
\[
y=(y^1,\dots,y^m,0,\dots,0),
\]
por lo que

\[
c(\pi(y))=c(y^1,\dots,y^m)=(y^1,\dots,y^m,0,\dots,0)=y.
\]
\begin{enumerate}[start=6]
\item Entonces
\end{enumerate}
\[
c\circ\pi=\operatorname{id}_{c(\mathbb R^m)}.
\]
\begin{enumerate}[start=7]
\item Así, $\pi=c^{-1}$ sobre la imagen, y $c:\mathbb R^m\to c(\mathbb R^m)$ es un homeomorfismo.
\end{enumerate}
\begin{itemize}
\item \textbf{Conclusión} La aplicación $c$ es suave, inyectiva, su diferencial es inyectivo en todo punto, y además $c$ es un homeomorfismo sobre su imagen. Por lo tanto
\end{itemize}
\[
c:\mathbb R^m\to\mathbb R^n
\]
es una incrustación suave.

\bookqed
\end{bookproof}
\begin{bookexample}{Helice}
Sea

\[
f:\mathbb R\to\mathbb R^3,\qquad f(t)=(\cos t,\sin t,t).
\]
Queremos ver que $f$ es una incrustación suave de $\mathbb R$ en $\mathbb R^3$.

\end{bookexample}
\begin{bookproof}{}
\begin{itemize}
\item \textbf{$f$ es suave}
\end{itemize}
\begin{enumerate}
\item Usando la carta de identidad de ambos lados con $\mathbb{R}$ y $\mathbb{R}^{3}$ respectivamente
\item Se reduce a ver que $f$ es suave como lo haciamos normalmente
\item Las funciones coordenadas de $f$ son
\end{enumerate}
\[
f^1(t)=\cos t,\qquad f^2(t)=\sin t,\qquad f^3(t)=t.
\]
\begin{enumerate}[start=4]
\item Todas son $C^\infty$, por lo tanto $f$ es $C^\infty$.
\end{enumerate}
\begin{itemize}
\item \textbf{$f$ es inyectiva}
\end{itemize}
\begin{enumerate}
\item Supongamos que
\end{enumerate}
\[
f(t)=f(s).
\]
\begin{enumerate}[start=2]
\item Entonces
\end{enumerate}
\[
(\cos t,\sin t,t)=(\cos s,\sin s,s).
\]
\begin{enumerate}[start=3]
\item Comparando la tercera coordenada,
\end{enumerate}
\[
t=s.
\]
\begin{enumerate}[start=4]
\item Por lo tanto $f$ es inyectiva.
\end{enumerate}
\begin{itemize}
\item \textbf{$df_t$ es inyectiva}
\end{itemize}
\begin{enumerate}
\item $df_t:T_t\mathbb R\to T_{f(t)}\mathbb R^3$ con base de $T_t\mathbb R$
\end{enumerate}
\[
\left\{\frac{\partial}{\partial t}\bigg|_t\right\},
\]
y base de $T_{f(t)}\mathbb R^3$:

\[
\left\{\frac{\partial}{\partial x_k}\bigg|_{f(t)}\right\}.
\]
\begin{enumerate}[start=2]
\item Luego
\end{enumerate}
\[
f'(t)=(df)_t\frac{\partial}{\partial t}\bigg|_t=\sum_{k=1}^3 a_k(t)\frac{\partial}{\partial x_k}\bigg|_{f(t)}.
\]
\begin{enumerate}[start=3]
\item Encontramos $a_k$ actuando sobre $x_k=\pi_k$ (aca asumimos que la carta para $\mathbb{R}^{3}$ es la identidad)
\end{enumerate}
\[
\sum_{k=1}^3 a_k(t)\frac{\partial}{\partial x_k}\bigg|_{f(t)}=\sum_{k=1}^3 (df)_t\frac{\partial}{\partial t}\bigg|_t(\pi_k)\frac{\partial}{\partial x_k}\bigg|_{f(t)}.
\]
\begin{enumerate}[start=4]
\item Por definición de diferencial:
\end{enumerate}
\[
\sum_{k=1}^3\frac{\partial}{\partial t}\bigg|_t(\pi_k\circ f)\frac{\partial}{\partial x_k}\bigg|_{f(t)}.
\]
\begin{enumerate}[start=5]
\item Ahora:
\end{enumerate}
\[
\pi_1\circ f=\cos t,\qquad \pi_2\circ f=\sin t,\qquad \pi_3\circ f=t.
\]
\begin{enumerate}[start=6]
\item Entonces:
\end{enumerate}
\[
(df)_t\left(\frac{\partial}{\partial t}\bigg|_t\right)=-\sin t\frac{\partial}{\partial x_1}\bigg|_{f(t)}+\cos t\frac{\partial}{\partial x_2}\bigg|_{f(t)}+1.\frac{\partial}{\partial x_3}\bigg|_{f(t)}
\]
\begin{enumerate}[start=7]
\item Identifico $T_{f(t)}\mathbb R^3$ con $\mathbb R^3.$ Entonces tenemos el viejo y conocido
\end{enumerate}
\[
(df)_t\left(\frac{\partial}{\partial t}\bigg|_t\right)=(-\sin t,\cos t,1).
\]
\begin{enumerate}[start=8]
\item Luego
\end{enumerate}
\[
(df)_t\lambda\frac{\partial}{\partial t}\bigg|_t
\]
es cero sii $\lambda=0,$ porque

\[
\lambda(-\sin t,\cos t,1)=0.
\]
\begin{enumerate}[start=9]
\item Por lo tanto
\end{enumerate}
\[
\ker(df_t)=\{0\}.
\]
Osea $df_{t}$ es inyectiva para cualquier $t$

\begin{itemize}
\item \textbf{$f$ es homeomorfismo sobre su imagen}
\end{itemize}
\begin{enumerate}
\item La imagen de $f$ es
\end{enumerate}
\[
f(\mathbb R)=\{(\cos t,\sin t,t):t\in\mathbb R\}\subseteq\mathbb R^3.
\]
\begin{enumerate}[start=2]
\item Definimos
\end{enumerate}
\[
\pi_3:\mathbb R^3\to\mathbb R,\qquad \pi_3(x^1,x^2,x^3)=x^3.
\]
\begin{enumerate}[start=3]
\item Entonces $\pi_3$ es continua. Además,
\end{enumerate}
\[
\pi_3(f(t))=\pi_3(\cos t,\sin t,t)=t.
\]
\begin{enumerate}[start=4]
\item Por lo tanto
\end{enumerate}
\[
\pi_3|_{f(\mathbb R)}:f(\mathbb R)\to\mathbb R
\]
es la inversa de $f:\mathbb R\to f(\mathbb R)$. 5. En efecto:

\[
\pi_3|_{f(\mathbb R)}\circ f=\operatorname{id}_{\mathbb R},
\]
y si $p\in f(\mathbb R)$, entonces $p=f(t)$ para algún $t$, y

\[
f(\pi_3(p))=f(\pi_3(f(t)))=f(t)=p.
\]
\begin{enumerate}[start=6]
\item Entonces
\end{enumerate}
\[
f\circ \pi_3|_{f(\mathbb R)}=\operatorname{id}_{f(\mathbb R)}.
\]
\begin{enumerate}[start=7]
\item Así, $f$ es un homeomorfismo sobre su imagen.
\end{enumerate}
\begin{itemize}
\item \textbf{Conclusión.} La aplicación $f$ es suave, inyectiva, su diferencial es inyectivo en todo punto, y además $f$ es un homeomorfismo sobre su imagen. Por lo tanto
\end{itemize}
\[
f:\mathbb R\to\mathbb R^3
\]
es una incrustación suave.

\bookqed
\end{bookproof}
\begin{bookexercise}{}
En general, si

\[
f:I\subseteq\mathbb R\to M
\]
es suave, entonces $f$ es inmersión si y solo si

\[
f'(t)\neq 0
\]
para todo $t\in I$.

\end{bookexercise}
\begin{bookproof}{}
\begin{enumerate}
\item Para cada $t\in I$, el espacio tangente $T_tI$ es $1$-dimensional y tiene base
\end{enumerate}
\[
\left\{\frac{\partial}{\partial t}\bigg|_t\right\}.
\]
\begin{enumerate}[start=2]
\item Por definición, la derivada de la curva en $t$ es
\end{enumerate}
\[
f'(t):=(df)_t\left(\frac{\partial}{\partial t}\bigg|_t\right)\in T_{f(t)}M.
\]
\begin{enumerate}[start=3]
\item Supongamos primero que $f$ es inmersión. Entonces
\end{enumerate}
\[
(df)_t:T_tI\to T_{f(t)}M
\]
es inyectiva para todo $t\in I$. 4. Como $\frac{\partial}{\partial t}\big|_t\neq 0$, por inyectividad se tiene

\[
(df)_t\left(\frac{\partial}{\partial t}\bigg|_t\right)\neq 0.
\]
Por lo tanto $f'(t)\neq 0$ para todo $t\in I$. 5. Recíprocamente, supongamos que $f'(t)\neq 0$ para todo $t\in I$. 6. Sea $v\in T_tI$. Como $T_tI$ es $1$-dimensional, existe $\lambda\in\mathbb R$ tal que

\[
v=\lambda\frac{\partial}{\partial t}\bigg|_t.
\]
\begin{enumerate}[start=7]
\item Si $(df)_t(v)=0$, entonces
\end{enumerate}
\[
0=(df)_t\left(\lambda\frac{\partial}{\partial t}\bigg|_t\right)=\lambda (df)_t\left(\frac{\partial}{\partial t}\bigg|_t\right)=\lambda f'(t).
\]
\begin{enumerate}[start=8]
\item Como $f'(t)\neq 0$, necesariamente $\lambda=0$. Luego $v=0$.
\item Entonces
\end{enumerate}
\[
\ker(df)_t=\{0\},
\]
y por lo tanto $(df)_t$ es inyectiva para todo $t\in I$. 10. Concluimos que $f$ es inmersión si y solo si $f'(t)\neq 0$ para todo $t\in I$.

\bookqed
\end{bookproof}
\begin{bookexample}{Esfera es subvariedad incrustada}
Sea $S^2$ la esfera con su estructura de variedad suave usual, y pensemos en

\[
i:S^2\to\mathbb R^3
\]
la inclusión,

\[
p\mapsto i(p)=p.
\]
Veamos que $(S^2,i)$ es una subvariedad incrustada.

\end{bookexample}
\begin{bookproof}{}
\begin{itemize}
\item \textbf{$i$ es suave}
\end{itemize}
\begin{enumerate}
\item Tomamos la carta identidad en $\mathbb{R}^{3}$ y la proyeccion estereografica $\varphi$ (o cualquiera otra) y se reduce a ver que dado $p\in S^{2}$ $i\circ\varphi ^{-1}$ es suave
\item Pero $i\circ\varphi ^{-1}=\varphi ^{-1}$ por definicion de $i$ y recordemos que $\varphi ^{-1}$ es suave por ser $\varphi$ carta
\end{enumerate}
\begin{itemize}
\item \textbf{$i$ es función inyectiva} Trivial
\item \textbf{$(di)_p:T_pS^2\to T_p\mathbb R^3$ es inyectiva}
\end{itemize}
\begin{enumerate}
\item Supongamos que $p$ está en el casquete $U_1^\varepsilon$ y Tenemos la carta
\end{enumerate}
\[
(U_{i}^{\epsilon,} \psi_1^\varepsilon=(x_1,x_2)),
\]
dada por $q=(q_{1},q_{2},q_{3})$

\[
x_1(q)=q_2,\qquad x_2(q)=q_3.
\]
\begin{enumerate}[start=2]
\item Su inversa es
\end{enumerate}
\[
(\psi_1^\varepsilon)^{-1}(s_1,s_2)=(\varepsilon\sqrt{1-s_1^2-s_2^2},s_1,s_2).
\]
\begin{enumerate}[start=3]
\item Recordar:
\end{enumerate}
\[
U_i^\varepsilon=\{(q_1,q_2,q_3)\in S^2:q_i\varepsilon>0\},\qquad \varepsilon=\pm1.
\]
\begin{enumerate}[start=4]
\item Tenemos base de $T_pS^2$ dada por
\end{enumerate}
\[
\mathcal B_1=\left\{\frac{\partial}{\partial x_1}\bigg|_p,\frac{\partial}{\partial x_2}\bigg|_p\right\}.
\]
\begin{enumerate}[start=5]
\item Y base de $T_p\mathbb R^3$ dada por
\end{enumerate}
\[
\mathcal B_2=\left\{\frac{\partial}{\partial y_1}\bigg|_p,\frac{\partial}{\partial y_2}\bigg|_p,\frac{\partial}{\partial y_3}\bigg|_p\right\}.
\]
\begin{enumerate}[start=6]
\item Calculamos la matriz de $[(di)_p]^{\mathcal B_2}_{\mathcal B_1}.$
\end{enumerate}
\[
(di)_p\frac{\partial}{\partial x_1}\bigg|_p=\sum_{k=1}^3 a_k(p)\frac{\partial}{\partial y_k}\bigg|_p.
\]
\begin{enumerate}[start=7]
\item Aplicamos $y_k$ para obtener $a_k$:
\end{enumerate}
\[
\sum_{k=1}^3 (di)_p\frac{\partial}{\partial x_1}\bigg|_p(y_k)\frac{\partial}{\partial y_k}\bigg|_p.
\]
\begin{enumerate}[start=8]
\item Por definición de diferencial:
\end{enumerate}
\[
\sum_{k=1}^3 \frac{\partial}{\partial x_1}\bigg|_p(y_k\circ i)\frac{\partial}{\partial y_k}\bigg|_p=\sum_{k=1}^3 \frac{\partial}{\partial s_1}\bigg|_{\varphi_{1}^{\epsilon} (p)}(y_k\circ i\circ(\psi_{1}^{\epsilon})^{-1} )\frac{\partial}{\partial y_k}\bigg|_p.
\]
\begin{enumerate}[start=9]
\item Ademas notemos que
\end{enumerate}
\[
y_1\circ i\circ(\psi_1^\varepsilon)^{-1}(s_1,s_2)=\varepsilon\sqrt{1-s_1^2-s_2^2},
\]
\[
y_2\circ i\circ(\psi_1^\varepsilon)^{-1}(s_1,s_2)=s_1,
\]
\[
y_3\circ i\circ(\psi_1^\varepsilon)^{-1}(s_1,s_2)=s_2.
\]
\begin{enumerate}[start=10]
\item Y ademas notando que $\psi_1^\varepsilon(p)=(p_{2},p_{3})$ tenemos
\end{enumerate}
\[
(di)_p\frac{\partial}{\partial x_1}\bigg|_p=\varepsilon\frac{-p_2}{\sqrt{1-p_2^2-p_3^2}}\frac{\partial}{\partial y_1}\bigg|_p+1\frac{\partial}{\partial y_2}\bigg|_p+0\frac{\partial}{\partial y_3}\bigg|_p.
\]
\begin{enumerate}[start=11]
\item Similarmente:
\end{enumerate}
\[
(di)_p\frac{\partial}{\partial x_2}\bigg|_p=\varepsilon\frac{-p_3}{\sqrt{1-p_2^2-p_3^2}}\frac{\partial}{\partial y_1}\bigg|_p+0\frac{\partial}{\partial y_2}\bigg|_p+1\frac{\partial}{\partial y_3}\bigg|_p.
\]
\begin{enumerate}[start=12]
\item Por lo tanto la matriz de $[(di)_p]^{\mathcal B_2}_{\mathcal B_1}$ es
\end{enumerate}
\[
\begin{pmatrix}-\frac{p_2}{p_1}&-\frac{p_3}{p_1}\\1&0\\0&1\end{pmatrix}.
\]
\begin{enumerate}[start=13]
\item Por el teorema del rango y la dimensión, $(di)_p$ es inyectiva sii $\operatorname{Im}(di)_p$ tiene rango $2$.
\item Y esto es equivalente a que la anterior matriz tenga una submatriz $2\times 2$ no singular.
\item En este caso sería $A[2:3,2:3]$ con
\end{enumerate}
\[
A=[(di)_p]^{\mathcal B_2}_{\mathcal B_1}.
\]
\begin{itemize}
\item \textbf{Homeomorfismo}
\end{itemize}
\begin{enumerate}
\item Es trivial, por que la topologia que le dimos a $S^{2}$ para verla como variedad diferenciable es la topologia heredada de $\mathbb{R}^{3}$
\end{enumerate}
\bookqed
\end{bookproof}
\begin{bookexample}{El ocho no es subvariedad incrustada}
Considérese

\[
f:\mathbb{R}\to\mathbb R^2, \qquad f(t)=\sin t\left(\cos t,1\right)=\left(\frac{\sin 2t}{2},\sin t\right)
\]
Su imagen es la curva con forma de ocho. Entonces $((0,2\pi),f)$ es una subvariedad de $\mathbb{R}^{2}$ pero no incrustada

\end{bookexample}
\begin{bookproof}{}
2. $g$ es claramente suave e inyectiva 3. Devuelta usando cartas identidad y bases tangentes podemos calcular el jacobiano (gradiente en este caso).

\[
[dg_{t}]=\begin{pmatrix}\cos^{2}t-\sin^{2}t\\\cos t\end{pmatrix}=\begin{pmatrix}\cos2t\\ \cos t\end{pmatrix}
\]
\begin{enumerate}[start=4]
\item Y esta claro que la unica forma de que esto sea el vector nulo es que $\cos t=0$ osea $t=\frac{2k+1}{2}\pi$ entonces $\cos 2t=\cos(2k+1)=1$. Osea no puede ser el vector nulo. Para ningun punto $t$
\item Por ende la unica forma de que el diferencial sea $0$ es evaluando en $0$ osea su $\ker = \{ 0 \}$
\item Pero $f$ no es inyectiva asi que no es subvariedad
\item Entonces redefinimos a
\end{enumerate}
\[
g:(0,2\pi)\to\mathbb R^2, \qquad g(t)=\sin t\left(\cos t,1\right)
\]

\bookimage{assets/pasted-image-20260425112115.png}{Pasted image 20260425112115}


\begin{itemize}
\item \textbf{$f$ es inyectiva} Sea $f(t)=\left(\frac{\sin 2t}{2},\sin t\right)=\left(\frac{\sin 2s}{2},\sin s\right)=f(s)$.
\end{itemize}
\begin{enumerate}
\item Si $t=s$, listo. Por lo que supongamos $t\neq s$.
\item De $\sin t=\sin s$ se sigue que $t=s+2k\pi$ o $t=\pi-s+2k\pi$, con $k\in\mathbb Z$.
\item Si $\sin t=0$, entonces $\sin s=0$, y en $(0,2\pi)$ solo es posible si $t=s=\pi$.
\item Si $\sin t\neq0$, entonces $\sin s\neq0$. Además, de $\frac{\sin 2t}{2}=\frac{\sin 2s}{2}$ se obtiene $\sin t\cos t=\sin s\cos s$.
\item Como $\sin t=\sin s\neq0$, concluimos $\cos t=\cos s$.
\item En cualquiera de los casos, $\sin t=\sin s$ y $\cos t=\cos s$.
\item Como $0<s,t<2\pi$, esto implica $|t-s|<2\pi$.
\item Pero si tienen el mismo seno y coseno, entonces $t-s=2k\pi$ con $k\in\mathbb Z$.
\item Luego $|k|<1$, por lo tanto $k=0$. Así, $t=s$.
\end{enumerate}
\begin{itemize}
\item Que es inmersion sale igual que antes por que el diferencial no cambia. Luego $((0,2\pi),g)$ es subvariedad.
\item \textbf{$f$ no es homeomorfismo}
\end{itemize}
\begin{enumerate}
\item Para ver esto, tenemos que mostrar que $g$ no produce un homeomorfismo entre $(0,2\pi)$ y el ocho con la topología de $\mathbb R^2$.
\item La idea está en que hay puntos "cercanos" en el ocho pero cuando se ven en $(0,2\pi)$ están "lejos".
\item Razonemos por el absurdo. Supongamos que
\end{enumerate}
\[
g:(0,2\pi)\to\text{ocho}
\]
es homeomorfismo. 4. Entonces

\[
g^{-1}:\text{ocho}\to(0,2\pi)
\]
es continua. 5. Tomemos la sucesión de puntos del ocho

\[
p_n=g\left(\frac1n\right)=\left(\frac{\sin(2/n)}{2},\sin(1/n)\right).
\]
\begin{enumerate}[start=6]
\item Con la topología del subespacio, $\{p_n\}$ es una sucesión de puntos del ocho que converge a $(0,0)$.
\item Por tanto, como $g^{-1}$ es continua, la sucesión $\{g^{-1}(p_n)\}$ debería converger a $g^{-1}(0,0)=\pi$.
\item Pero
\end{enumerate}
\[
g^{-1}(p_n)=\frac1n,
\]
luego

\[
g^{-1}(p_n)\to0.
\]
\begin{enumerate}[start=9]
\item Por lo tanto, $g$ no es homeomorfismo sobre su imagen.
\end{enumerate}
\begin{itemize}
\item Así, $((0,2\pi),g)$ es una subvariedad inmersa, pero no una subvariedad incrustada.
\end{itemize}
\bookqed
\end{bookproof}
\section{Equivalencia de subvariedades}
\begin{bookdefinition}{Equivalencia de subvariedades}
Dos subvariedades

\[
(N_1,f_1),\qquad (N_2,f_2)
\]
de una variedad suave $M$ se dicen \textbf{equivalentes} si existe un difeomorfismo

\[
\Phi:N_1\to N_2
\]
tal que

\[
f_1=f_2\circ\Phi
\]
Esto define una relación de equivalencia en el conjunto de todas las subvariedades de $M$. Analogamente 
\bookimage{assets/pasted-image-20260425120519.png}{Pasted image 20260425120519}
 es un diagrama conmutativo

\end{bookdefinition}
\begin{bookproposition}{Representante incrustado único}
Sea $M$ una variedad suave. Toda clase de equivalencia $[(N,f)]$ de subvariedades de $M$ contiene un único representante de la forma

\[
(\widehat N,i)
\]
donde $\widehat N\subseteq M$ es un subconjunto e

\[
i:\widehat N\hookrightarrow M
\]
es la inclusión.

\end{bookproposition}
\begin{bookproof}{}
\begin{enumerate}
\item Como $(N,f)$ una subvariedad de $M$. Como $f$ es inyectiva, induce una biyección
\end{enumerate}
\[
\widehat f:N\to \widehat N:=f(N)
\]
notar que como cambia el conjunto de llegada cambia la funcion 2. Transportamos la topología y la estructura suave de $N$ a $\widehat N$ mediante $\widehat f$. Con esa estructura, $\widehat f$ es un difeomorfismo. 3. Si $i:\widehat N\hookrightarrow M$ es la inclusión, entonces

\[
f=i\circ \widehat f
\]
\begin{enumerate}[start=4]
\item Veamos que $i:\widehat N\to M$ es suave. Las cartas de $\widehat N$ son de la forma $(\widehat U,\psi=\varphi\circ \widehat f^{-1})$, donde $\widehat U=\widehat f(U)$ y $(U,\varphi)$ es carta de $N$.
\item Para cada una de estas cartas,
\end{enumerate}
\[
i\circ(\varphi\circ \widehat f^{-1})^{-1}:\varphi(U)\to M
\]
es suave, pues

\[
i\circ(\varphi\circ \widehat f^{-1})^{-1}= f\circ\varphi^{-1}:\varphi(U)\to M,
\]
que es suave, ya que $f:U\to M$ es suave por hipótesis y $\varphi$ es carta, osea suave 6. Notar que alcanza con una sola de las composiciones por que si $f\circ \varphi ^{-1}$ es suave entonces $\psi\circ f\circ\varphi ^{-1}$ va a ser suave para cualquier $\psi$ carta\\
7. En cuanto a $(di)_p:T_p\widehat N\to T_pM$, esta es inyectiva, por que $f=i\circ\widehat f$, entonces $df_p=(di)_{\widehat f(p)}(d\widehat f)_p$. 8. Y al ser $(d f)_p$ inyectiva y $(d\widehat f)_{p}$ isomorfismo sobre su imagen (por que $\tilde{f}$ es homeomorfismo), entonces $(di)_{\widehat f(p)}$ debe ser inyectiva. 9. Por lo tanto $(\widehat{N},i)$ es claramente subvariedad y por 3. (y 2. que dice $\widehat{f}$ es difeo) tenemos $(\widehat{N},i)\in [(N,f)]$

\begin{itemize}
\item ¿Es el único representante en $[N,f]$ de esta forma?
\end{itemize}
\begin{enumerate}
\item Si $(\widetilde N,i)\in[N,f]$, esto quiere decir que existe
\end{enumerate}
\[
g:N\to\widetilde N
\]
difeo tal que

\[
f=i\circ g.
\]
\begin{enumerate}[start=2]
\item Pero $f=i\circ\widehat f$  por lo que $\widetilde N=\widehat N$ como conjuntos y  $g=\widehat f$.
\item Ahora queremos ver que $\widehat{N}$ y $\widetilde{N}$ tienen la misma estructura diferenciable $\widehat{\mathcal{F}}$ y $\widetilde{\mathcal{F}}$. (Sabemos que son iguales como conjuntos)
\item Miramos
\end{enumerate}
\[
\operatorname{Id}_{\widehat N}:(\widetilde N,\widetilde{\mathcal F})\to(\widehat N,\widehat{\mathcal F}).
\]
\begin{enumerate}[start=5]
\item Como $\widehat f:N\to(\widehat N,\widehat{\mathcal A})$ es difeo y $\widehat f=g:N\to(\widetilde N,\widetilde{\mathcal A})$ también es difeo (por ser inmersiones, definicion), entonces
\end{enumerate}
\[
g\circ\widehat f^{-1}=\widehat f\circ\widehat f^{-1}=\operatorname{Id}_{\widehat N}
\]
es difeo. (Que la $Id_{\widehat{N}}$ sea difeo nos dice que para cualquier par de cartas $(\widetilde{N},\varphi)$ y $(\widehat{N},\psi)$ sucede $\psi\circ Id_{\widehat{N}}\circ\varphi ^{-1}=\psi\circ\varphi ^{-1}$ es suave. Mostrando la compatibilidad) 6. Por ende las estructuras son las mismas, por que si tomamos cartas cualquier $\varphi,\psi$ en $\widetilde{F},\widehat{F}$ respectivamente

\bookqed
\end{bookproof}
\section{Consecuencia práctica}
\begin{bookremark}{Identificar una subvariedad con su imagen}
Gracias a la proposición anterior, muchas veces una subvariedad se identifica con un subconjunto

\[
\widehat N\subseteq M
\]
provisto de la estructura suave que hace a la inclusión una incrustación.

En ese contexto, decir que “$\widehat N$ es una subvariedad de $M$” significa precisamente que la inclusión

\[
i:\widehat N\hookrightarrow M
\]
es una incrustación suave.

\end{bookremark}
\section{Ejemplo final}
\begin{bookexample}{Dos presentaciones del ocho}
Considérese la imagen del ocho en $\mathbb R^2$ y las dos presentaciones

\[
f:(0,2\pi)\to\mathbb R^2,\qquad f|_{(0,2\pi)}(t)=\left(\frac{\sin 2t}{2},\sin t\right),
\]
y

\[
\tilde f:(-\pi,\pi)\to\mathbb R^2,\qquad \tilde f|_{(-\pi,\pi)}(t)=\left(\frac{\sin 2t}{2},\sin t\right)
\]
Estas dos subvariedades no son equivalentes. 
\bookimage{assets/pasted-image-20260425170149.png}{Pasted image 20260425170149}


\begin{bookremark}{}
\begin{enumerate}
\item Si fueran equivalentes, existiría un difeomorfismo
\end{enumerate}
\[
F:(0,2\pi)\to(-\pi,\pi)
\]
tal que

\[
f=\tilde f\circ F
\]
\begin{enumerate}[start=2]
\item La unica $F$ que cumple eso manda $t$ a $t$ para $0<t<\pi$ $t$ a $t-2t$ para $\pi<t<2\pi$ y $t$ a $\pi$ si $t=\pi$
\item \bookimage{assets/pasted-image-20260425170729.png}{Pasted image 20260425170729}
\item ¿Cómo sale $F$? Si
\end{enumerate}
\[
f|_{(-\pi,\pi)}\circ F=f|_{(0,2\pi)}.
\]
\begin{enumerate}[start=5]
\item Si $t\in(0,\pi)$, entonces $f|_{(-\pi,\pi)}(t)=f|_{(0,2\pi)}(t)$.
\item Por tanto,
\end{enumerate}
\[
f|_{(-\pi,\pi)}(F(t))=f|_{(0,2\pi)}(t)=f|_{(-\pi,\pi)}(t).
\]
\begin{enumerate}[start=7]
\item Y como $f|_{(-\pi,\pi)}$ es inyectiva, se sigue que $F(t)=t$.
\item Si $t=\pi$, entonces
\end{enumerate}
\[
f|_{(-\pi,\pi)}(F(\pi))=f|_{(0,2\pi)}(\pi)=(0,0).
\]
\begin{enumerate}[start=9]
\item Y como $f|_{(-\pi,\pi)}(0)=0$, con $f|_{(-\pi,\pi)}$ inyectiva, se sigue que
\end{enumerate}
\[
F(\pi)=0.
\]
\begin{enumerate}[start=10]
\item Ya con esto, $F$ no es continua. Ver el tramo $(\pi,2\pi)$.
\end{enumerate}
\end{bookremark}
\end{bookexample}
\phantomsection\label{blk:8ed2bd}
